\appendix
\chapter{Registro de cambios} % usa \section si es article
\begin{longtable}{|p{3.5cm}|p{5.5cm}|p{5.5cm}|}
\hline
\textbf{Sección} & \textbf{Comentario} & \textbf{Cambio realizado} \\
\hline
\endfirsthead
\hline
\textbf{Sección} & \textbf{Comentario} & \textbf{Cambio realizado} \\
\hline
\endhead
\hline
\endfoot
1. Introducción & Reestructurar de lo general a lo particular; primero hablar de non-profits, que es el capítulo específico del SCN donde entran las CAC (Sebastián Cea) & Se agregaron dos párrafos iniciales: uno sobre el sector sin fines de lucro y la economía social en el marco del SCN (Manual ONU-TSE 2018), y otro sobre las cooperativas en general, incluyendo los distintos tipos investigados (trabajo asociado, agrícolas, consumo, vivienda, servicios y financieras). El párrafo sobre CAC se reformuló como transición a lo particular. \\
\hline
Introducción / Marco institucional & Corregir el nombre institucional: se menciona ``División de Asociatividad (DAES)'' cuando la institución que efectivamente apoyó el proyecto Huella Social es el Instituto Nacional de Asociatividad (Sebastián Cea) & Se eliminó del cuerpo de la memoria (Sección 1.1, Marco institucional del proyecto) el párrafo que listaba a la "División de Asociatividad y Cooperativismo (DAES)" como socio estratégico del proyecto. Se reemplazó por una remisión directa al Anexo~\ref{anx:fondef}, donde se corrigió el listado de instituciones según el formulario FONDEF ID26 (ver fila "Anexos -- Formulario FONDEF"). El cuerpo del texto ahora identifica correctamente al Instituto Nacional de Asociatividad y Cooperativismo (INAC) como socio institucional del proyecto. \\
\hline
Introducción & Unificar el nombre de la División de Asociatividad y Cooperativismo (DAES), que aparecía escrito de dos formas distintas en el capítulo ("Cooperativismo" en un párrafo y "Cooperativas" en otro) & Se unificó el nombre a "División de Asociatividad y Cooperativismo (DAES)" en ambas menciones del capítulo, consistente con la denominación usada en el formulario FONDEF ID26 y en el informe de evaluación del concurso. \\
\hline
\hline
Resultados -- Cuenta de producción total & Hallazgo propio detectado en la revisión (no proviene literalmente de la reunión): el texto afirmaba en dos lugares (intro del Capítulo~\ref{ch:resultados} y Sección~\ref{sec:limitaciones}, punto 3) que era "metodológicamente improcedente" sumar los segmentos CMF y DAES, pero la Sección~\ref{sec:agregado-total} sí realiza y reporta esa suma, generando una contradicción interna & Se reformuló la intro del Capítulo~\ref{ch:resultados} y el punto 3 de la Sección~\ref{sec:limitaciones} para aclarar que el agregado CMF+DAES es una aproximación de orden de magnitud (no una suma metodológicamente exacta), en vez de calificarlo como "improcedente". Se corrigió además el período de DAES citado en Limitaciones (decía 2015--2024; se corrigió a 2014--2024, consistente con el resto del documento). Se agregó una oración de enganche en la Sección~\ref{subsec:validacion-cmf} recordando esta misma salvedad. \\
\hline
\hline
Resultados -- Validación PIB & Realizar una validación cruzada entre los resultados de aporte al PIB estimados y la información publicada por la Comisión para el Mercado Financiero (CMF); si los resultados difieren, documentar los supuestos y posibles razones (Jose Joaquín Fernández / Antonio Ruiz-Tagle). Sebastián Cea aportó el criterio específico de validación: si el aporte al PIB del sistema bancario es de aprox. 5\,\%, los resultados de las CAC deberían ser coherentes con su proporción patrimonial respecto de la banca (2,4\,\%) & Se agregó la Sección~\ref{subsec:validacion-cmf} ("Validación con evidencia externa") al final de la Sección~\ref{sec:agregado-total} (Agregado total: CMF y DAES). Se verificó el 2,4\,\% mencionado en la reunión contra una fuente oficial (Informe de Desempeño CMF, dic. 2025: patrimonio CAC = 2,41\,\% del patrimonio conjunto banca+cooperativas), y se sustituyó el "5\,\% de aporte bancario al PIB" (dato no verificable, mencionado solo de forma verbal en la reunión) por la participación de la actividad "Servicios financieros" en el PIB según el Banco Central de Chile (3,1\,\%--4,7\,\%, 2013--2025), con la salvedad metodológica de que esa categoría no equivale exactamente a "banca". Se construyó el Cuadro~\ref{tab:validacion_cmf}, comparando año a año el aporte esperado bajo este criterio contra el aporte efectivamente estimado (CMF+DAES). El orden de magnitud resultó consistente en la mayoría de los años (diferencia $<$0,01 pp en 2013, 2014 y 2021), con una brecha de hasta 0,03 pp entre 2017--2020 atribuible al supuesto de $\alpha$ uniforme. \\
\hline
\hline
Introducción -- Alcances (Sección 1.3) & Integrar el formulario FONDEF, actualmente en los anexos, con referencias explícitas dentro del cuerpo de la memoria, evitando que parezcan documentos desconectados y mostrando que el proyecto de título representa una mejora y evolución sobre el formulario original (Sebastián Cea) & Se reemplazó la oración final de la sección por un párrafo que presenta la delimitación del estudio como la estrategia de pilotaje definida en la postulación FONDEF: las CAC como caso piloto de la plataforma de cuentas satélite, con escalamiento posterior al conjunto de la ESS (remisión al Anexo~\ref{anx:fondef}). \\
\hline
Marco teórico -- Brechas del estado del arte (Sección 3.5.2) & Mismo comentario (integración del formulario FONDEF) & Se reformuló el párrafo final de la sección: se explicita que el diagnóstico de fragmentación de registros fundamenta la postulación FONDEF (Anexo~\ref{anx:fondef}) y que la memoria entrega sustento cuantitativo, para un clúster específico de organizaciones, al análisis de magnitud que el formulario planteaba en términos descriptivos. \\
\hline
Metodología -- Fuentes de información & Mismo comentario & Se agregó un párrafo al cierre del inventario de fuentes (Sección~\ref{sec:fuentes-informacion}) que lo vincula con las 15 fuentes de datos declaradas en el formulario, contrastando la amplitud de la postulación con la profundidad de la memoria (cobertura, método de extracción y brechas por fuente), y recogiendo la observación del panel evaluador sobre la ausencia de procesos explícitos de verificación y normalización de los datos. \\
\hline
Metodología -- Panel interactivo & Mismo comentario & Se agregó un párrafo de cierre en la Sección~\ref{sec:panel} que sitúa el panel como avance concreto de la hoja de ruta tecnológica comprometida en el formulario, en tanto eslabón intermedio entre el prototipo Shiny R previo al proyecto y la plataforma en React comprometida en la postulación. \\
\hline
Resultados -- Validación con evidencia externa & Mismo comentario & Se agregó una oración en la Sección~\ref{subsec:validacion-cmf} que explicita que la validación cruzada implementa una práctica cuya ausencia señaló el panel evaluador de la postulación, dejando establecido un procedimiento replicable para el desarrollo futuro de la plataforma. \\
\hline
Conclusiones (Sección 6.2) & Mismo comentario & Se reformularon los párrafos previo y posterior a la lista de debilidades del panel evaluador: se agregó el folio de la postulación (ID26I10768) y se precisó que la memoria aborda dos de las cuatro debilidades (la cuantificación del problema por tipo de organización, ejecutada para el caso CAC, y la validación cruzada de la Sección~\ref{subsec:validacion-cmf}), explicitando la evolución verificable de la memoria respecto del formulario original. \\
\hline
Anexos -- Formulario FONDEF & Mismo comentario & Se podaron redundancias en el Anexo A.1.1 mediante remisiones a las Secciones 3.3 y 3.5.2 del cuerpo, y se agregó al cierre del Anexo A.1.2 un párrafo de correspondencia que mapea los tres objetivos específicos del formulario a las secciones donde la memoria los ejecuta (Sección~\ref{sec:fuentes-informacion}; Sección~\ref{sec:metodos-estimacion}; Secciones~\ref{sec:panel} y~\ref{sec:dashboard}). \\
\hline
\hline
Análisis y conclusiones & Incluir en el análisis y conclusiones una revisión de los comentarios recibidos en la evaluación FONDEF (Folio ID26I10768, nota final 2,73/5, no aprobado), como guía estratégica para orientar la ejecución de futuras memorias (Sebastián Cea) & Se agregó un párrafo al final de la Sección "Conclusiones" (antes de "Recomendaciones para trabajos futuros") que resume las cuatro debilidades principales señaladas por el panel evaluador de FONDEF, indica cuál de ellas fue parcialmente abordada en este trabajo (la validación cruzada de la Sección~\ref{subsec:validacion-cmf}) y conecta las restantes con la sección de recomendaciones. \\
\hline
1.2 Objetivos específicos & Añadir texto explicativo que conecte los objetivos específicos entre sí y con el capítulo 1.3, indicando en qué parte del documento se desarrolla cada objetivo (Sebastián Cea) & Se agregó un párrafo en la Sección 1.2.2 (Objetivos específicos), justo después del listado, que explica la secuencia lógica entre los cuatro objetivos y referencia mediante los capítulos y secciones donde se desarrolla cada uno: objetivo 1 en los Capítulos~\ref{ch:revision} y~\ref{ch:marco}; objetivo 2 en la Sección~\ref{sec:fuentes-informacion} del Capítulo~\ref{ch:metodologia}; objetivo 3 en la Sección~\ref{sec:metodos-estimacion} del Capítulo~\ref{ch:metodologia} y en el Capítulo~\ref{ch:resultados}; y objetivo 4 en la Sección~\ref{sec:panel} del Capítulo~\ref{ch:metodologia} y la Sección~\ref{sec:dashboard} del Capítulo~\ref{ch:resultados}. \\
\hline
Resultados -- Gráficos del tablero interactivo & Ajustar las imágenes de los gráficos del tablero interactivo, ya que se veían cortadas en los bordes; explorar su inserción directa mediante código Python en Overleaf en vez de capturas de pantalla (Jose Joaquín Fernández / Sebastián Cea) & Se generó el código en Python que produce los gráficos directamente (sin recurrir a capturas de pantalla), eliminando el recorte en los bordes y generando las figuras completas para su inserción en el documento. \\
\hline
\end{longtable}